%%%%%%%%%%%%%%%%%%%%%%%%%%%%%%%%%%%%%%%%%%%%%%%%%%%%%%%%%%%%%%%%%%%%%%%%%%%%%%
% BYU / Utah NASA Space Grant Consortium Fellowship -- Proposal Narrative
% Applicant: Sam Charles (undergraduate)
% Cycle:     2026 (deadline: 8 May 2026)
% Length:    3-5 pages (per program guidelines, see nasa-26-application.pdf)
%
% First-draft narrative populated from Sam's outline (see PR comment thread on
% issue #26).  Assumptions and gaps that need Sam's input are flagged with the
% \needsinput{...} macro -- search the source (or look for the bracketed
% [needs input: ...] callouts in the rendered PDF) to find every spot that
% requires a decision before submission.  Open questions are also collected
% in the final "Open questions" section so they are easy to scan.
%
% Build (MiKTeX):
%   cd proposals/byu-nasa-space-grant-2026
%   latexmk -pdf proposal.tex
%
% Required MiKTeX packages (auto-installed when [MPM]AutoInstall=1) -- keep
% this list in sync with the \usepackage lines below:
%   geometry, times, microtype, graphicx, titlesec, fancyhdr, lastpage,
%   xcolor, soul, enumitem, natbib, hyperref
%%%%%%%%%%%%%%%%%%%%%%%%%%%%%%%%%%%%%%%%%%%%%%%%%%%%%%%%%%%%%%%%%%%%%%%%%%%%%%

\documentclass[11pt]{article}

\usepackage[letterpaper,margin=1in]{geometry}
\usepackage{times}
\usepackage{microtype}
\usepackage{graphicx}
\usepackage{titlesec}
\usepackage{fancyhdr}
\usepackage{lastpage}
\usepackage{xcolor}
\usepackage{soul}
\usepackage{enumitem}
\usepackage[round,authoryear]{natbib}
\usepackage[hidelinks]{hyperref}

% Inline callout for assumptions / gaps that Sam still needs to confirm.
% Renders as a highlighted "[needs input: ...]" so it is impossible to miss
% in the PDF, and is also easy to grep for in the source.
\sethlcolor{yellow}
\newcommand{\needsinput}[1]{\hl{[needs input: #1]}}

% Tighter, bolder section headings appropriate for a short proposal narrative.
\titleformat{\section}{\normalfont\large\bfseries}{\thesection}{0.6em}{}
\titleformat{\subsection}{\normalfont\normalsize\bfseries}{\thesubsection}{0.6em}{}
\titlespacing*{\section}{0pt}{1.0\baselineskip}{0.3\baselineskip}
\titlespacing*{\subsection}{0pt}{0.6\baselineskip}{0.2\baselineskip}

\setlength{\parskip}{0.4\baselineskip}
\setlength{\parindent}{0pt}

\pagestyle{fancy}
\fancyhf{}
\fancyhead[L]{\small Sam Charles --- NASA Space Grant Fellowship Proposal}
\fancyhead[R]{\small Page \thepage\ of \pageref{LastPage}}
\renewcommand{\headrulewidth}{0.4pt}

\begin{document}

\begin{center}
  {\Large\bfseries Generative CAD for an Open-Source Multi-Powder Doser:\\
   A Tooling Step Toward Autonomous Discovery of Additively Manufactured
   Aerospace Alloys}\\[6pt]
  {\large Sam Charles\\
   Undergraduate Applicant, Brigham Young University}\\[2pt]
  {\small Utah NASA Space Grant Consortium Fellowship --- 2026 Cycle}
\end{center}
\vspace{0.4\baselineskip}

\section{Introduction and Motivation}

Additively manufactured (AM) structural alloys are an increasingly important
class of materials for space applications, where launch-mass reduction,
topology-optimized geometries, and rapid iteration on bespoke parts all
favor laser powder-bed fusion (L-PBF) over conventional subtractive
processes \citep{placeholder_am_aerospace}.  Discovering and qualifying
\emph{new} aerospace alloys for these processes -- rather than only printing
established compositions -- is currently bottlenecked not by the printers
themselves, but by the upstream workflow that prepares the metal powder
feedstock from which each candidate composition is built.

A central, repetitive step in that workflow is \emph{powder dosing}:
weighing out and combining individual elemental or pre-alloyed powders into
a target blend before each L-PBF build or screening experiment.  In a
research lab this is almost always done by hand, one powder at a time, with
a balance and a spatula.  This is slow, error-prone, exposes researchers to
fine reactive metal powders, and -- critically -- it caps the
\emph{compositional throughput} of any downstream alloy-discovery
campaign.  The few commercial automated dispensers that do exist are
priced for industrial QC labs (typically tens of thousands of dollars),
designed around pharmaceutical or analytical-chemistry workflows rather
than L-PBF feedstock preparation, and ill-matched to the
volume\,/\,particle-size regime that aerospace alloy printing actually
needs \needsinput{confirm with Sam: typical price quotes and model numbers
of the commercial dispensers we are positioning against -- pending
literature/landscape sweep in issue~\#28~/~\#10}.

Powder dispensing is therefore an \emph{unmet need} in the research
laboratory setting.  Existing solutions -- including the open-source
designs we have surveyed -- fall short of the rigorous design and testing
required for an L-PBF feedstock workflow, where the dispenser must
respect aerospace-grade powder characteristics (e.g.\ tightly controlled
particle-size distribution, sphericity, oxygen pickup) and must deliver
total dispense volumes large enough to actually fill a build hopper
\citep{placeholder_powder_dosing}.

This project will close that gap.  We will use a mix of traditional
mechanical design and \emph{generative CAD} -- driven primarily by
agentic AI tooling layered on top of GitHub Copilot -- with frequent
build-test-feedback cycles between AI-generated designs and the human
engineers (Sam and Will) to create a powder dosing system capable of
handling \textbf{up to 30 unique powders} with \textbf{total dispense
volumes up to 250\,mL}, suitable as the front-end of an autonomous
aerospace-alloy discovery loop.

The work has two complementary novel contributions:
\begin{enumerate}[leftmargin=1.3em,itemsep=0.1em,topsep=0.2em]
  \item A rigorous, documented exploration of the \emph{capabilities and
        current limitations} of agentic AI applied to CAD and
        electromechanical design, using a real, non-trivial aerospace-relevant
        device as the test article.
  \item An \emph{open-source}, accurate, low-cost, reliable multi-powder
        dispenser that is purpose-built for L-PBF research workflows --
        released as a complete repository (CAD, firmware, BOM, build
        instructions, and test data) so other research groups can adopt
        and extend it.
\end{enumerate}

A peer-reviewed paper documenting both contributions is planned as a
follow-on deliverable.

\section{Background}

\subsection{Powder dispensing in research and industry}

Accurate dispensing of fine metal powders at research scale is dominated
today by a small number of commercial systems aimed at analytical-chemistry
and pharmaceutical applications.  These systems achieve high single-dose
accuracy, but at price points that are out of reach for most academic
materials labs and with form factors and software stacks that assume a
single feedstock at a time -- exactly the opposite of what an alloy
combinatorial campaign needs \needsinput{cite specific commercial systems
once the issue~\#28 / \#10 landscape sweep produces a reference list; the
\texttt{placeholder\_commercial\_dispensers} entry in \texttt{refs.bib}
is a stand-in until then}.

On the open-source / DIY side, prior work has demonstrated single-powder
augers, vibratory dispensers, and small carousel feeders, but to our
knowledge none of these published designs has been validated end-to-end
against L-PBF feedstock requirements -- particle-size-distribution
preservation, cross-contamination between powders, oxygen exposure, and
bulk volume per dispense -- in the way an aerospace alloy discovery
workflow demands \citep{placeholder_powder_dosing}
\needsinput{any specific open-source projects you want explicitly named
and cited?  e.g.\ a particular GitHub repo or thesis; pending the
\#28 background sweep}.

\subsection{Generative CAD as a design accelerator}

\emph{Generative CAD} here refers to the broader family of techniques in
which a designer specifies intent (functional requirements, constraints,
interfaces) and software produces candidate geometry, rather than the
designer drawing the geometry directly.  In this project the ``generator''
is an agentic large-language-model loop -- primarily GitHub Copilot
backed by multiple model providers, augmented with task-specific tooling
for CAD scripting, slicing, and printer control (cf.\ companion issues
\#6 and \#24 in the project repository) -- not a topology-optimization
solver or a parametric template library, although those are complementary
techniques the loop is free to invoke
\citep{placeholder_generative_cad,placeholder_agentic_ai_engineering}.

The published research base on this style of LLM-driven, end-to-end
electromechanical design is still thin: most of the literature is
restricted either to topology optimization for a single fixed envelope
or to short, demonstrative ``can the model produce a bracket'' studies.
The open question is whether a tightly-supervised agentic loop can
deliver an \emph{integrated} mechanical + electrical + firmware design
for a real research instrument.  This proposal directly probes that
question \needsinput{add concrete generative-CAD references once
issue~\#28 produces them; until then the bib entry
\texttt{placeholder\_generative\_cad} is a stand-in}.

\subsection{Connection to autonomous AM alloy discovery}

Autonomous, ``self-driving'' materials laboratories have shown the value
of closing experimental loops with active learning over composition,
process, and structure-property data \citep{placeholder_self_driving_lab}.
Applied to AM aerospace alloys, the loop is gated by how quickly and
reproducibly we can prepare new powder blends.  An accurate, scriptable,
30-powder dispenser is the missing front-end for that loop and the
direct technical contribution of this fellowship.

\section{Project Objectives}

\begin{enumerate}[leftmargin=1.5em,itemsep=0.15em,topsep=0.2em]
  \item \textbf{Open-source multi-powder doser (hardware).}  A working
        prototype that handles up to 30 distinct powders with per-blend
        total dispense volumes up to 250\,mL, with documented dispense
        accuracy, repeatability, and cross-contamination characterization
        on at least \needsinput{target number, e.g.\ 3?} representative
        L-PBF feedstock powders.
  \item \textbf{Open-source repository.}  A complete release at
        \href{https://github.com/vertical-cloud-lab/powder-doser}%
             {\texttt{github.com/vertical-cloud-lab/powder-doser}}
        \citep{placeholder_powder_doser_repo} containing CAD sources,
        firmware, BOM, assembly instructions, test data, and the AI
        prompts/transcripts that produced each design revision.
  \item \textbf{Generative-CAD-first design process.}  As little
        traditional, hand-driven CAD as we can manage -- ideally the
        engineers should never need to open a CAD GUI to author or edit
        geometry, only to inspect and review what the agentic loop
        produced.  Every deviation from that ideal will be logged and
        analyzed.
  \item \textbf{Documented capability/limitation study.}  A thorough,
        reproducible record of where agentic AI succeeded, where it
        failed, where engineer intervention was required, and what
        tooling unlocked each capability.  This deliverable functions as
        a litmus test for the current state of generative CAD applied to
        a real electromechanical instrument and will be the basis of the
        follow-on paper.
\end{enumerate}

It is important to state up front what this project is \emph{not}: this
is not an evaluation of AI's ability to develop a product unsupervised.
Capable engineers (Sam and Will) will review the AI's contributions at
every step and make meaningful corrections.  The fellowship work is an
honest evaluation of agentic AI as a \emph{tool} that augments capable
engineers, with generative CAD as the sharpest test case.

\section{Approach and Methods}

\subsection{Generative CAD pipeline}

The design loop is anchored on GitHub Copilot used in agent mode, with
the freedom to call out to multiple underlying language models and to
auxiliary tooling exposed through the project's repository (e.g.\ the
component-selection work in \#24 for the vibration motor / solenoid
control stack on a Raspberry Pi Zero 2~W, and the meta-CAD tooling
survey in \#6).  Each design iteration follows a build-test-feedback
cycle:
\begin{enumerate}[leftmargin=1.5em,itemsep=0.1em,topsep=0.2em]
  \item The engineers describe a requirement or failure mode in plain
        language and reference data (test results, photos of failures,
        powder-flow observations).
  \item The agentic loop produces or edits CAD sources, electrical
        layouts, and firmware -- preferring scriptable, version-controlled
        formats so the diff is human-reviewable.
  \item Prototypes are printed in-house on a Bambu Lab H2D
        \needsinput{confirm: H2D is the printer we will use end-to-end?
        Any second printer (resin? metal?) we should mention?}.  The
        long-term goal is for the agent to drive the printer
        \emph{directly} using the H2D programmatic-print pathway being
        developed in \#23, so the loop can iterate without an
        engineer-in-the-loop print step; in the near term the engineers
        will trigger prints manually.
  \item Sam and Will assemble, bench-test, and report results back into
        the loop.  Each iteration's prompts, model outputs, and engineer
        corrections are committed to the repo so the design history is
        fully auditable.
\end{enumerate}

\subsection{Hardware design targets}

The dispenser must, at minimum, satisfy the following functional
requirements derived from the L-PBF feedstock-preparation use case:
\needsinput{Sam: please confirm/adjust each of the targets below before
submission -- I have inferred reasonable values from the project context
but they have not been formally agreed.}
\begin{itemize}[leftmargin=1.3em,itemsep=0.1em,topsep=0.2em]
  \item Up to 30 distinct powder reservoirs, individually addressable.
  \item Total per-blend dispense volume up to 250\,mL.
  \item Per-powder dispense resolution and accuracy compatible with
        compositional screening at the
        \needsinput{target accuracy: e.g.\ $\pm 1$\,wt\%? $\pm 0.1$\,wt\%?}
        level.
  \item Particle-size-distribution preservation across the dispense
        path (no significant attrition or segregation of the feedstock).
  \item Negligible cross-contamination between adjacent powder
        reservoirs.
  \item Compatibility with reactive aerospace alloy powders
        \needsinput{do we need an inert-atmosphere enclosure as part of
        the deliverable, or is this out of scope for the fellowship
        year?}.
\end{itemize}

\subsection{Validation and integration with the autonomous discovery loop}

Validation will combine bench-level tests of the dispenser in isolation
(gravimetric repeatability over hundreds of doses, PSD measurement
before/after, cross-contamination checks) with an integration test in
which the dispenser prepares blends consumed by a downstream L-PBF or
small-scale alloy-screening experiment
\needsinput{which downstream platform will consume the blends for the
integration test? a specific L-PBF machine, an arc-melt screening
setup, or both?}.  Results feed directly back into the generative-CAD
loop both as design feedback and as the dataset that the follow-on paper
will report.

\section{Timeline and Deliverables}

The bulk of the engineering work is targeted for the summer (May--August
2026), when both engineers (Sam and Will) have the most concentrated
availability.  The remainder of the calendar year is reserved for
characterization, documentation, and writing.

\begin{itemize}[leftmargin=1.3em,itemsep=0.1em,topsep=0.2em]
  \item \textbf{May 2026 (Months 0--1):} Finalize requirements; complete
        the issue~\#28 background sweep; stand up the agentic CAD loop
        end-to-end against a stub problem; select and procure mechanical
        + electrical components (cf.\ \#24).
  \item \textbf{June 2026 (Month 2):} First full-system prototype
        (\(\le\)\,4 reservoir subset) printed on the H2D, basic
        dispense-accuracy bench tests, first round of generative-CAD
        capability/limitation observations logged.
  \item \textbf{July 2026 (Month 3):} Scale-up to the full 30-reservoir
        envelope; close the loop on direct printer control via \#23 so
        the agent can iterate prints without engineer intervention;
        cross-contamination and PSD-preservation tests.
  \item \textbf{August 2026 (Month 4) -- soft completion:} Integrated
        dispenser delivered, repository public-ready, capability/limitation
        study draft complete.
  \item \textbf{Fall 2026 (Months 5--7):} Integration test with the
        downstream alloy-screening workflow; expanded characterization;
        first paper draft.
  \item \textbf{Winter 2027 (Months 8--9):} Paper revisions and
        submission; community release of v1.0 of the open-source
        repository.
\end{itemize}

\section{Broader Impact and Relevance to NASA}

The most direct line to NASA is through aerospace structural alloys: a
research-grade, low-cost, open-source multi-powder dispenser shortens
the compositional-discovery loop for the L-PBF alloys that go into
launch vehicles, in-space structures, and propulsion components.  The
same dispensing problem also appears, with different chemistries, in
several other NASA-relevant areas: precise dispensing of fine powders
for solid and hybrid rocket propellant formulation studies, regolith
simulant handling for in-situ resource utilization (ISRU) experiments,
and lunar/Martian regolith additive manufacturing research
\needsinput{Sam, do you want to lean harder on any specific NASA
mission directorate or program (e.g.\ STMD, Artemis, MSFC AM efforts)?
that hook would strengthen this section}.

The agentic-AI portion of the work is also directly relevant to NASA's
broader interest in AI-assisted engineering for cost- and
schedule-constrained missions: a documented, honest assessment of where
generative CAD currently helps and where it currently fails is exactly
the kind of evidence that mission engineering teams need before they can
adopt these tools.  Finally, releasing the dispenser as a fully
open-source design lowers the barrier of entry for other Utah and
NASA-affiliated labs that want to run AM alloy-discovery campaigns
without committing to industrial-scale capital equipment.

\section{Qualifications of the Applicant}

\needsinput{Sam to write -- short paragraph covering relevant coursework,
prior project experience (especially anything mechanical, electrical,
firmware, AM, or AI/CAD), and role on this project.  Mention Will and
his role here as well.  Also list faculty advisor and the cost-share
arrangement when known.}

\section*{Open Questions for the Applicant}

The first-draft prose above makes several reasonable-but-unconfirmed
assumptions.  Each is also marked inline with \texttt{[needs input: \dots]}
so they are easy to find in the rendered PDF.  Consolidated:
\begin{enumerate}[leftmargin=1.5em,itemsep=0.1em,topsep=0.2em]
  \item Confirm or adjust the hardware targets in
        \S\,Approach (per-powder accuracy, need for inert atmosphere,
        number of validation powders).
  \item Confirm the printer(s) we are committing to in this proposal --
        is the Bambu H2D the only printer in scope, or do we also want
        to mention any second machine?
  \item Confirm the downstream consumer for the integration test (which
        L-PBF or alloy-screening platform will eat the blends).
  \item Decide which specific commercial dispensers, open-source prior
        art, and generative-CAD references we want explicitly cited;
        these are the gaps the issue~\#28 / \#10 background sweeps are
        meant to fill.
  \item Decide whether to lean on a particular NASA program/center hook
        in the Broader Impact section.
  \item Provide the Qualifications paragraph for Sam (and Will, where
        relevant), plus faculty advisor and cost-share details.
\end{enumerate}

\textit{Suggested additions to consider:} a one-figure schematic of the
dispenser concept (would help the narrative considerably -- a hand sketch
or an early CAD render is fine); a brief letter-of-support style mention
of Will's role; and a short sentence on data/code-availability
commitments (open license, archived release, etc.).

\section*{References}
\vspace{-1.0\baselineskip}
\renewcommand{\bibsection}{}
\bibliographystyle{plainnat}
\bibliography{refs}

\end{document}
